% -*- mode : latex; coding: utf-8 -*-
\begin{KoreanAbstract}
  딥러닝, 인공지능 등 현대의 데이터센터 워크로드들은 큰 메모리 용량을 요구한다. 그러나, 프로세서에 메모리를 꽂을 수 있는 핀의 개수는 제한되어 있다. 이 메모리 부족 문제를 해결하지 못한다면 미래의 데이터 센터 성능에 저하가 생길 수 있다. 가능한 해결 방법 중 하나는 CXL \cite{CXL}을 이용한 메모리 분산이다. CXL로 연결된 메모리를 사용함으로써 프로세서에 있는 메모리 핀 개수와 무관하게 메로리 용량을 늘릴 수 있다. 본 연구에서는 트레이스를 기반으로 사이클 레벨에서 CXL 타입 3 장치의 성능을 측정할 수 있는 CXLmem 시뮬레이터를 구현했다. 이 시뮬레이터를 트레이스를 기반으로 사이클 레벨에서 x86 구조의 성능을 측정할 수 있는 MacSim 시뮬레이터 붙여서, SPEC CPU 2006 워크로드 중 메모리 부하가 큰 워크로드 8개로 4 코어 시뮬레이션을 진행했다. 그 결과, CXL로 연결된 메모리 장치는 로컬 DRAM 메모리만 쓸 때에 비해 평균 20\%, 최대 40\%의 IPC (사이클 당 명령어 개수) 하강을 유발했다. 또한 메모리에서 데이터를 읽는데 걸리는 시간의 경우, CXL로 연결된 메모리 장치가 로컬 DRAM 메모리에 비해 평균 2배 긴 시간을 보였다. CXL로 연결된 메모리를 로컬 DRAM 메모리와 함께 사용할 경우에는 로컬 DRAM에 가해지는 부하가 작아져, 로컬 DRAM 메모리에서 데이터를 읽는데 걸리는 시간이 18\% 감소했다.
\end{KoreanAbstract}
